% !TEX encoding = UTF-8 Unicode

\documentclass[aspectratio=169]{beamer}
\usepackage[utf8]{inputenc}

%% BIB
%\usepackage[
%  style=authoryear,
%  backend=biber,
%  url=false,
%  maxcitenames=2,
%  uniquename=false,
%  uniquelist=false
%]{biblatex}
%\addbibresource{?.bib}

%% GRAPHICS

\usepackage{graphicx}
\graphicspath{{../images/},{../figures/}}
\usepackage{tikz}

%% COLORS
\definecolor{UWRed}{HTML}{C5050C}
\definecolor{TintedBG}{HTML}{EEEEFF}
%\definecolor{StrongBlue}{HTML}{3F8FD2}
%\definecolor{StrongGreen}{HTML}{36C88E}
%\definecolor{StrongRed}{HTML}{9B0000}
%\definecolor{MyC}{HTML}{009999}
%\definecolor{MyM}{HTML}{990099}
%\definecolor{MyY}{HTML}{999900}
%\definecolor{MyR}{HTML}{990000}
%\definecolor{MyG}{HTML}{009900}
%\definecolor{MyB}{HTML}{000099}
%\definecolor{ActionRed}{HTML}{990000}

%% SLIDE COLOR SETTINGS
\setbeamercolor{structure}{fg=UWRed}
%\setbeamercolor{title page}{fg=white}
%\setbeamercolor{title}{fg=white}

%% RM NAV SYMBOLS
\setbeamertemplate{navigation symbols}{}

%% FONTS
\setbeamerfont{title}{size=\huge\bfseries}

%% DRAWING
%\usetikzlibrary{}

%% LOGO on slides
\logo{\begin{tikzpicture}[overlay]
  \node[anchor=north east,inner sep=0] at (0,86mm) {\includegraphics[height=10mm]{SMPH_color-flush.pdf}};
\end{tikzpicture}}

%% CONTENT BEGINS

\title{Analysis of high-throughput mass spectrometry of CRISPR knockout cells}
\subtitle{Craven Lab Group Meeting}
\author{Yuriy Sverchkov}
\institute{University of Wisconsin--Madison}
\date{November 16, 2018}

\begin{document}

  {
    %\setbeamertemplate{background canvas}{\includegraphics[width=\paperwidth]{UW-lake.png}}
    \begin{frame}[plain]
      %\vskip2cm
      \titlepage
    \end{frame}
  }

%%% FRAMEBREAK %%%

\begin{frame}{The H3K Project}
\begin{columns}
	\column{0.4\textwidth}
\begin{itemize}
	\item Goal: Study mitochondrial function
	\item Mitochondrial dysfunction is associated with various diseases of interest
	\item However, many of the molecular processes involved are poorly characterized
	\item Approach: Multi-omic measurements of CRISPR knockout cell lines
\end{itemize}
\column{0.57\textwidth}
\includegraphics[width=\textwidth]{h3k-knockouts}
\end{columns}
\end{frame}

%%% FRAMEBREAK %%%

\begin{frame}[plain]
\includegraphics[width=\linewidth]{h3k-tech}
\end{frame}

%%% FRAMEBREAK %%%

\begin{frame}[plain]
\includegraphics[width=\linewidth]{h3k-molecules}
\end{frame}

%%% FRAMEBREAK %%%

\begin{frame}{The data}
\begin{itemize}
	\item Currently working with $K=29$ (sentinel) knockouts and $M=13080$ measured biomolecules
	\item We are interested in the relative effect of a KO on the measurements
	\item Analyses in the following is all based on log-2-fold-change of molecule abundance in KO with respect to wild type
	\item $x_{mk} =$ $t$-statistic of $\log_2$-fold-change in the abundance of molecule $m$ in KO $k$ with respect to wild type, taken over replicates
	\item[] $\left( \frac{\text{mean}}{\text{stdev}\times \sqrt{\text{\# replicates}} } \right)$
\end{itemize}
\end{frame}

%%% FRAMEBREAK %%%

\begin{frame}{Building a correlational network}
\begin{columns}
	\column{0.5\textwidth} \small
	Similarity measure between molecules $i$ and $j$
	\begin{itemize}
		\item Euclidean Distance:
		\[
		\sqrt{ \sum_{k} (x_{ik} - x_{jk})^2 }
		\]\pause
		\item Pearson Correlation:
		\[
		\frac{\sum_k ( x_{ik} - \bar{x_k} ) ( x_{jk} - \bar{ x_k } ) }{ \sqrt{ \sum_k ( x_{ik} - \bar{x_k} )^2 } \sqrt{ \sum_k ( x_{jk} - \bar{ x_k } )^2 } }
		\]\pause
		\item Cosine Similarity: \alert{$\leftarrow$ our choice}
		\[
		\frac{\sum_k x_{ik} x_{jk}}{ \sqrt{ \sum_k x_{ki\strut}^2 } \sqrt{ \sum_k x_{jk\strut}^2 } }
		\]\pause
	\end{itemize}
	\column{0.47\textwidth} \small
	Thresholding rule to obtain a network from pairwise cosine similarities based on a null distribution:
	\pause
	
	\begin{itemize}
	\item $t \sim$ cosine similarity of two random (i.i.d. normal) vectors in $\mathbb R^K$
	\item $(t + 1)/2 \sim \mathrm{Beta}( \frac{K - 1}{2}, \frac{K - 1}{2} )$ \pause
	\item Using a Bonferroni correction $\alpha = 0.5/M$
	\item Keep $s_{ij}:$
	\[ | s_{ij} | > 1 - 2 \mathrm{cdf}^{-1}_{\mathrm{Beta}\left( \frac{K - 1}{2}, \frac{K - 1}{2} \right) }\left(  \alpha \right) \]
	\end{itemize}
\end{columns}
\end{frame}

%%% FRAMEBREAK %%%

\begin{frame}{Cosine similarity network}
\begin{columns}
	\column{0.7\textwidth}
	\includegraphics[width=\textwidth]{full-similarity-network.png}
	\column{0.27\textwidth}
	\begin{tabular}{l|r}
		\multicolumn{2}{c}{Connected Nodes} \\ \hline
		Proteins & 2130 \\
		Metabolites & 133 \\
		Lipids & 1742 \\ \hline
		Total & 4005 \\ \hline \hline
		\multicolumn{2}{c}{Edges} \\ \hline
		P-P & 9585 \\
		P-M & 397 \\
		P-L & 2743 \\
		M-M & 183 \\
		M-L & 572 \\
		L-L & 11113 \\ \hline
		Total & 24593
	\end{tabular}
\end{columns}
\end{frame}

%%% FRAMEBREAK %%%

\begin{frame}{Cosine similarity network annotation}
Checking the network against known interaction databases by ``annotating'' it, i.e., matching network edges to documented interactions

\begin{itemize}
	\item STRING Consortium protein-protein interaction database
	\item KEGG interactions:
	\begin{itemize}
		\item \textbf{Compound} (Metabolite or Lipid) \textbf{- Compound} interactions: do the two compounds participate in the same reaction
		\item \textbf{Protein - Compound} interactions: does the protein map to an enzyme that participates in a reaction with the compound
		\item \textbf{Protein - Protein} interactions: do the two proteins share a pathway
	\end{itemize}
\end{itemize}

\end{frame}

%%% FRAMEBREAK %%%

\begin{frame}{Annotated similarity network statistics}
\centering
\begin{tabular}{c|r|crr}
	\multicolumn{2}{c}{Similarity Edges} & STRING & KEGG interaction & KEGG ID collision \\ \hline
	P-P & 9585 & 100 (23 overlap) & 185 & 1 \\
	P-M & 397 & & 0 & \\
	P-L & 2743 & & 0 & \\
	M-M & 183 & & 2 & 3 \\
	M-L & 572 & & 0 & \\
	L-L & 11113 & & 200 & 547 \\ \hline
	Total & 24593 & 100 & 387
\end{tabular}

\pause
\vfill
\begin{itemize}
	\item A very small fraction of the similarity edges could be matched to background knowledge
	\item Mappings between our molecules and databases are not one-to-one
	\item Particularly lipids don't map well to KEGG IDs
\end{itemize}
\end{frame}

%%% FRAMEBREAK %%%

\begin{frame}{Annotated similarity network}
\centering
\only<1>{\includegraphics[width=7cm]{annotated-similarity-network}}
\only<2>{\includegraphics[width=7cm]{annotated-parts-only}}
\end{frame}

%%% FRAMEBREAK %%%

\begin{frame}{Building a nested effects model}
\centering
\includegraphics[width=0.5\textwidth]{fgnem-title}
%\hfill
%\includegraphics[width=0.2\textwidth]{plos-logo}


\includegraphics[width=0.5\textwidth]{fgnem}
\end{frame}

%%% FRAMEBREAD %%%
\begin{frame}{Nested effect model}
\centering
\includegraphics[width=9cm]{fgnem-graph}
\end{frame}

\end{document}